\section{Introduction 引言}

本文围绕研究问题展开，首先说明其应用场景以及现实意义。开头需要明确交代当前方法的局限，并阐明开展这项工作的实践动机或科学动机。行文应尽量使用具体、清晰的语言，同时避免提前展开属于后续章节的实现细节。

本文所要解决的核心空缺可以概括为以下方面：现有方法在质量、效率、鲁棒性、可扩展性、可解释性、数据可得性或部署约束中的一个或多个方面仍存在不足。这里应以能够自然引出所提出方法的方式描述这一空缺。

本文随后提出相应的方法，并在技术章节开始前，为读者建立对系统、算法、数据集或分析流程的整体认识。必要时可结合主要概览图进行说明，例如图~\ref{fig:teaser}。

本文的主要贡献如下：
\begin{itemize}
    \item \textbf{贡献一。} 用一到两句话说明第一项技术贡献或实证贡献。
    \item \textbf{贡献二。} 说明第二项贡献，并强调其新颖性或实用价值。
    \item \textbf{贡献三。} 说明在验证、数据集、基准、应用或分析方面的贡献。
\end{itemize}

最后，简要总结支撑本文论断的证据，包括评测设置、重要指标和总体结论；在最终数值结果确定并准备发布之前，这里不列出具体数字。

\newpage
